\documentclass[11pt,reqno]{amsart}

\usepackage{amsmath,amssymb,amsthm}
\usepackage[margin=1.15in]{geometry}
\usepackage[colorlinks=true,linkcolor=blue,citecolor=blue]{hyperref}

\newtheorem{theorem}{Theorem}[section]
\newtheorem{proposition}[theorem]{Proposition}
\newtheorem{lemma}[theorem]{Lemma}
\newtheorem{corollary}[theorem]{Corollary}
\newtheorem{conjecture}[theorem]{Conjecture}
\theoremstyle{definition}
\newtheorem{definition}[theorem]{Definition}
\newtheorem{question}[theorem]{Question}
\newtheorem{remark}[theorem]{Remark}

\newcommand{\HH}{\mathcal{H}}
\newcommand{\PP}{\mathcal{P}}
\newcommand{\BH}{\mathcal{B}(\HH)}
\newcommand{\N}{\mathbb{N}}
\newcommand{\R}{\mathbb{R}}
\newcommand{\U}{\mathcal{U}}
\newcommand{\pluf}{\pi}
\newcommand{\ZW}{\mathcal{Z}}
\newcommand{\supp}{\operatorname{supp}}
\newcommand{\cod}{\operatorname{codim}}

\title[Intimate subspaces and block filters]
{Intimate subspaces, block filters,\\ and the diagonalization of maximal projection
filters}

\author{David Feldman}
\address{Department of Mathematics and Statistics, University of New Hampshire, Durham, NH}

\author{Alexander Wilce}
\address{Department of Mathematics, Susquehanna University, Selinsgrove, PA}

\date{\today}

\begin{document}

\begin{abstract}
Fix an orthonormal basis of a separable Hilbert space and call the closed spans of its
subsets \emph{blocks}. We prove a zero--cofinite dichotomy answering a question of Fred
Galvin: a closed subspace of infinite codimension misses some infinite block entirely,
so every closed subspace, restricted to a suitable infinite block, is either everything
but finitely much or nothing at all. Using it we show that the filter of closed subspaces
generated by the blocks of an ultrafilter $\U$ on $\N$, together with the
finite-codimension subspaces, is a \emph{maximal} filter in the projection lattice
(a \emph{pluf}, in the terminology of the companion paper) if and only if $\U$ is
selective. For non-selective $\U$ the block filter is properly extendable; we prove
that its extensions cannot be steered away from the witness subspace by any blocker
built from disjointly supported vectors, nor by any subspace meeting a single
piece-block in dimension two or more, and we conjecture that the witness subspace
belongs to \emph{every} maximalization --- a rigidity phenomenon with no Boolean
analogue. We then
develop the theory of \emph{intimate} subspaces: those meeting, for every two-coloring
of the basis, at least one of the two color blocks nontrivially. A construction of
Gowers, reproduced here with proof, shows intimacy strictly exceeds its obvious
sufficient conditions; our dichotomy shows that example has, in a precise sense, the
only possible shape. The pay-off is a characterization: a nonprincipal pluf is
diagonalizable over a basis if and only if every member is intimate for that basis.
Consequently non-diagonalizability is witnessed by a single member, the transfinite
machinery of the companion paper collapses to a one-witness reduction, and the question
of ZFC examples localizes to a concrete question about the continuous masa.
\end{abstract}

\maketitle

\section{Introduction}\label{sec:intro}

Throughout, $\HH$ is a separable infinite-dimensional Hilbert space with a fixed
orthonormal basis $(e_n)_{n\in\N}$; for $S\subseteq\N$ we write
$\HH_S=\overline{\operatorname{span}}\{e_n:n\in S\}$ and call such subspaces
\emph{blocks}. A \emph{pluf} is a maximal filter in the lattice $\PP(\HH)$ of closed
subspaces; the companion paper \cite{FWI} develops the general theory. Here we study the
interaction of plufs with a basis, organized around three results.

First, a dichotomy (\S\ref{sec:dichotomy}). Gowers attributed to Fred Galvin the
question: must a closed subspace meeting every infinite block have finite codimension?
Theorem~\ref{thm:A} answers it in the sharp form that a subspace of infinite
codimension misses some infinite block \emph{entirely}, completing a sketch of
Gowers.
Relativized (Corollary~\ref{cor:B}), this says the restriction of any closed subspace to
a suitable infinite sub-block is either cofinite-dimensional or zero --- no middle
ground survives passage to a sub-block.

Second, a characterization of maximality (\S\ref{sec:block}). An ultrafilter $\U$ on
$\N$ generates, together with the finite-codimension subspaces, a filter $\Phi(\U)$ in
$\PP(\HH)$. Theorem~\ref{thm:C}: $\Phi(\U)$ is a pluf if and only if $\U$ is selective.
The proof combines Corollary~\ref{cor:B} with Mathias's theorem on happy families;
necessity is witnessed by an explicit subspace $W$ built from a partition without
selectors. We then study the pluf extensions of $\Phi(\U)$ in the non-selective case
and prove two rigidity results (Proposition~\ref{prop:noblock},
Lemma~\ref{lem:thin}): the natural candidates for a member separating an extension
from $W$ all fail, for reasons that appear structural rather than accidental. We
conjecture that $W$ belongs to every pluf extension of $\Phi(\U)$
(Conjecture~\ref{conj:rigidity}) --- maximalization in the projection lattice, unlike
in a Boolean algebra, would then be canalized through specific transverse subspaces.

Third, intimacy and diagonalization (\S\S\ref{sec:intimate}--\ref{sec:diag}). Call $W$
\emph{intimate} (for the fixed basis) if for every $A\subseteq\N$, $W\cap\HH_A\neq0$ or
$W\cap\HH_{\N\setminus A}\neq0$. Subspaces of finite codimension in an infinite block,
and their enlargements, are intimate; Gowers's construction
(Theorem~\ref{thm:gowers}) shows that these do not exhaust intimacy. Corollary~\ref{cor:B} then shows the
example is not an accident of the construction: \emph{any} intimate subspace beyond the
trivial examples must have dense zero-set (every infinite set contains an infinite $S$
with $W\cap\HH_S=0$) while its zero-set contains no complementary pair
(Proposition~\ref{prop:shape}) --- exactly the shape Gowers built.

No characterization of intimacy is needed for the application: what plufs consume is
its \emph{failure}, one witness at a time.

\begin{theorem}[= Theorem~\ref{thm:D}]
A nonprincipal pluf $\pluf$ is diagonalizable over a basis $\mathcal B$ --- meaning
$\{S:\HH_S^{\mathcal B}\in\pluf\}$ is an ultrafilter --- if and only if every member of
$\pluf$ is $\mathcal B$-intimate.
\end{theorem}

Since non-intimacy of a single member survives every enlargement of a filter, the
diagonalizability of all extensions of a family is decided by the family itself
(Proposition~\ref{prop:reduction}). Theorem~\ref{thm:E} builds such a family under CH:
a non-diagonalizable nonprincipal pluf obtained without spectral theory, and without
carrying a maximalization through the recursion. Section~\ref{sec:zfc} assembles what is known toward removing CH,
including the observation that countable nested families can never witness the
reduction, and the identification of the continuous masa as the natural candidate.

\section{The zero--cofinite dichotomy}\label{sec:dichotomy}

\begin{theorem}[Galvin's question; completing a sketch of Gowers]\label{thm:A}
Let $W$ be a closed subspace of $\HH$ of infinite codimension. Then there is an infinite
$S\subseteq\N$ with $W\cap\HH_S=0$.
\end{theorem}

\begin{proof}
Let $U=W^\perp$, so $\dim U=\infty$. For each $N$, the subspace
$\{u\in U:u_n=0\text{ for }n\le N\}=\ker(P_{[1,N]}|_U)$ has codimension at most $N$ in
$U$, hence is nonzero. Recursively choose $N_0=0<n_1<N_1<n_2<N_2<\cdots$ and unit
vectors $u_j\in U$ as follows: $u_j$ vanishes on $[1,N_{j-1}]$; $n_j>N_{j-1}$ is a
coordinate with $\delta_j:=|u_j(n_j)|>0$; and $N_j>n_j$ is large enough that
$\bigl(\sum_{n>N_j}|u_j(n)|^2\bigr)^{1/2}<\delta_j2^{-j-2}$.

Set $S=\{n_1<n_2<\cdots\}$ and let $x\in W\cap\HH_S$. Fix $j$. For $i<j$ we have
$n_i\le N_i\le N_{j-1}$, so $u_j(n_i)=0$; hence
\[
0=\langle x,u_j\rangle
 =x_{n_j}\overline{u_j(n_j)}+\sum_{i>j}x_{n_i}\overline{u_j(n_i)}.
\]
For $i>j$, $n_i>N_{i-1}\ge N_j$, so Cauchy--Schwarz bounds the sum by
$\|x\|\,\delta_j2^{-j-2}$, giving $|x_{n_j}|\le2^{-j-2}\|x\|$. Then
$\|x\|^2=\sum_j|x_{n_j}|^2\le\|x\|^2\sum_{j\ge1}4^{-j-2}=\tfrac1{48}\|x\|^2$, so $x=0$.
\end{proof}

\begin{remark}
Gowers's sketch (row-echelon form for a spanning family of $W^\perp$, distinguished
columns with rapidly decreasing tails) is exactly the recursion above; the convergence
bookkeeping is the only addition. We have not found the statement in the literature; it
has the flavor of a Bessaga--Pe{\l}czy\'nski selection and may be known to Banach-space
specialists.
\end{remark}

\begin{corollary}[Relativized dichotomy]\label{cor:B}
For every closed $W$ and every infinite $S_0\subseteq\N$: either $W\cap\HH_{S_0}$ has
finite codimension in $\HH_{S_0}$, or there is an infinite $S\subseteq S_0$ with
$W\cap\HH_S=0$.
\end{corollary}

\begin{proof}
Apply Theorem~\ref{thm:A} inside $\HH_{S_0}$ to the subspace $W\cap\HH_{S_0}$.
\end{proof}

Thus the conceivable middle behaviors --- infinite dimension and infinite codimension at
once --- can always be destroyed by passing to a sub-block. Write
\[
\ZW_W=\{S\in[\N]^\infty : W\cap\HH_S=0\},\qquad
\mathcal C_W=\{S\in[\N]^\infty : \cod_{\HH_S}(W\cap\HH_S)<\infty\};
\]
both are closed under infinite subsets, and Corollary~\ref{cor:B} says
$\ZW_W\cup\mathcal C_W$ is dense in $([\N]^\infty,\subseteq)$.

\section{Block filters and selectivity}\label{sec:block}

\begin{definition}
For an ultrafilter $\U$ on $\N$, let $\Phi(\U)$ be the filter in $\PP(\HH)$ generated by
$\{\HH_S:S\in\U\}$ together with all closed subspaces of finite codimension.
\end{definition}

$\Phi(\U)$ is proper: a generating intersection has the form $\HH_S\cap V$ with $S\in\U$
and $\cod V<\infty$, which is infinite-dimensional. Two criteria are immediate:

\begin{lemma}\label{lem:criteria}
$(a)$ $W\in\Phi(\U)$ iff $\cod_{\HH_S}(W\cap\HH_S)<\infty$ for some $S\in\U$.
$(b)$ $W$ generates a proper filter with $\Phi(\U)$ iff $\dim(W\cap\HH_S)=\infty$ for
every $S\in\U$.
\end{lemma}

\begin{proof}
$(a)$ If $\cod_{\HH_S}(W\cap\HH_S)=d<\infty$, then $V=(W\cap\HH_S)\oplus\HH_{\N\setminus
S}$ has codimension $d$ and $\HH_S\cap V\subseteq W$. Conversely if
$\HH_S\cap V\subseteq W$ with $\cod V<\infty$ then
$\cod_{\HH_S}(W\cap\HH_S)\le\cod V$. $(b)$ $X$ meets every finite-codimension subspace
nontrivially iff $\dim X=\infty$; apply this to $X=W\cap\HH_S$ within the generating
intersections.
\end{proof}

Recall that $\U$ is \emph{selective} (Ramsey) iff every partition of $\N$ into sets not
in $\U$ admits a selector in $\U$, iff $\U$ is minimal in the Rudin--Keisler order, and
that Mathias's theorem gives: for selective $\U$ and analytic
$\mathcal A\subseteq[\N]^\infty$ there is $S\in\U$ with $[S]^\infty\subseteq\mathcal A$
or $[S]^\infty\cap\mathcal A=\emptyset$ \cite{Mathias}.

\begin{theorem}\label{thm:C}
$\Phi(\U)$ is a pluf if and only if $\U$ is selective.
\end{theorem}

\begin{proof}
\emph{Sufficiency.} Let $\U$ be selective and $W$ closed. The set
$\{S:W\cap\HH_S\neq0\}$ is analytic (it is the projection of the closed set
$\{(S,x):x\in W,\ \|x\|=1,\ \supp x\subseteq S\}$). By Mathias, some $S\in\U$ has
$[S]^\infty$ inside or disjoint from it. In the disjoint case $W\cap\HH_S=0$: the block
$\HH_S\in\Phi(\U)$ blocks $W$. In the inside case no infinite subset of $S$ lies in
$\ZW_W$, so Corollary~\ref{cor:B} forces $S\in\mathcal C_W$, and $W\in\Phi(\U)$ by
Lemma~\ref{lem:criteria}$(a)$. Every $W$ being decided, $\Phi(\U)$ is maximal (by the
maximality criterion of \cite{FWI}, Lemma 2.1: a filter is maximal iff every subspace is
a member or meets some member trivially).

\emph{Necessity.} Let $\U$ be non-selective and fix a partition
$\N=\bigsqcup_kA_k$ with no $A_k\in\U$ and no selector in $\U$. Put
$f_k=\sum_{n\in A_k}2^{-n}e_n$ and $W=\{f_k:k\}^\perp$. The restrictions
$P_Sf_k$, over $k$ with $A_k\cap S\neq\emptyset$, have pairwise disjoint supports, hence
are linearly independent; so for every $S$,
\[
\dim(W\cap\HH_S)=\sum_k\bigl(|A_k\cap S|-1\bigr)^+,\qquad
\cod_{\HH_S}(W\cap\HH_S)=\#\{k:A_k\cap S\neq\emptyset\}
\]
(with the convention $\infty-1=\infty$). For $S\in\U$: if the codimension were finite,
$S$ would lie in a finite union of the $A_k$ up to a finite set, putting some $A_k$ in
$\U$; so $W\notin\Phi(\U)$ by Lemma~\ref{lem:criteria}$(a)$. If the dimension were
finite, deleting finitely many points of $S$ would leave a set meeting each $A_k$ at
most once --- a selector in $\U$; so $\dim(W\cap\HH_S)=\infty$ for all $S\in\U$, and by
Lemma~\ref{lem:criteria}$(b)$, $W$ is addable. Addable and unadded: $\Phi(\U)$ is not
maximal.
\end{proof}

We now examine the extensions of $\Phi(\U)$ in the non-selective case. By the
maximality criterion, a pluf extension of $\Phi(\U)$ \emph{omits} $W$ if and only if
it has a member $R$ with $R\cap W=0$; and every member of every extension meets every
member of $\Phi(\U)$, i.e.\ satisfies $\dim(R\cap\HH_S)=\infty$ for all $S\in\U$
(Lemma~\ref{lem:criteria}$(b)$). Call such an $R$ an \emph{addable blocker} of $W$.
One might expect addable blockers to abound, as they do in the Boolean setting, where
a filter omitting a set always extends to an ultrafilter omitting it (adjoin the
complement). The next two results show that in $\PP(\HH)$ the natural candidates all
fail, in a way that suggests the failure is structural.

\begin{lemma}[Blockers are thin on every piece]\label{lem:thin}
If $R$ is closed with $R\cap W=0$, then $\dim\bigl(R\cap\ell^2(A_k)\bigr)\le1$ for
every $k$.
\end{lemma}

\begin{proof}
If $\dim(R\cap\ell^2(A_k))\ge2$, the single linear condition
$\langle\,\cdot\,,f_k\rangle=0$ cuts it in codimension at most one, leaving a nonzero
$v\in R\cap\ell^2(A_k)\cap f_k^\perp$; and $v\perp f_j$ for $j\neq k$ automatically
(disjoint supports), so $0\neq v\in R\cap W$.
\end{proof}

\begin{proposition}[No disjointly supported blocker]\label{prop:noblock}
Let $R$ be the closed span of a family $(\rho_i)$ of nonzero vectors with pairwise
disjoint supports $\sigma_i=\supp\rho_i$, and suppose $R\cap W=0$. Then there is
$S\in\U$ with $R\cap\HH_S=0$. In particular $R$ is not addable, and no pluf extension
of $\Phi(\U)$ omitting $W$ has a member of this form.
\end{proposition}

\begin{proof}
Since the supports are disjoint, the normalized $\rho_i$ are orthonormal, so every
$x\in R$ has a unique expansion $x=\sum_ic_i\rho_i$, and the restriction of $x$ to
$\sigma_i$ is $c_i\rho_i|_{\sigma_i}$. Hence $x\in\HH_S$ iff for every $i$ either
$c_i=0$ or $\sigma_i\subseteq S$:
\[
R\cap\HH_S=\overline{\operatorname{span}}\{\rho_i:\sigma_i\subseteq S\}.
\]
It therefore suffices to produce $S\in\U$ containing no $\sigma_i$, i.e.\ a set
$D\notin\U$ meeting every $\sigma_i$.

Let $D_s$ be the set of points carrying singleton supports. If two of them, $n\neq
n'$, lay in one piece $A_k$, then $e_n,e_{n'}\in R$ (scalar multiples of the two
generators) and the vector $2^{-n'}e_n-2^{-n}e_{n'}$ would lie in $R\cap W$: it is
orthogonal to $f_k$ by the choice of coefficients and to every other $f_j$ by
disjointness of supports. So $R\cap W=0$ forces $|D_s\cap A_k|\le1$ for all $k$:
$D_s$ is a selector of the partition, hence $D_s\notin\U$.

For each $i$ with $|\sigma_i|\ge2$, choose two distinct points
$p_i,q_i\in\sigma_i$. The sets $P=\{p_i\}$ and $Q=\{q_i\}$ are disjoint, so at most
one of them lies in the ultrafilter; let $D_1\in\{P,Q\}$ with $D_1\notin\U$. Then
$D=D_s\cup D_1$ is a union of two non-members, hence $D\notin\U$, and
$S=\N\setminus D\in\U$. Every $\sigma_i$ meets $D$ (singletons through $D_s$, the
rest through $D_1$), so no $\sigma_i\subseteq S$ and $R\cap\HH_S=0$.
\end{proof}

Lemma~\ref{lem:thin} says a blocker can carry at most a line's worth of each
piece-block; Proposition~\ref{prop:noblock} says that assembling such lines (or any
disjointly supported family, aligned with the pieces or not) never yields
addability, because the ultrafilter always contains the complement of a transversal
of the supports. Blockers, if they exist at all, must therefore couple infinitely
many pieces through vectors with overlapping supports, in such a way that
$R\cap\HH_S$ is sustained by cancellation rather than by whole generators. Three
further constraints on such an $R$ are proved in \cite{FWIV}, \S6: its intersection
with the block of any $m$ pieces has dimension at most $m$; its total support lies in
$\U$; and inside every $\HH_S$, $S\in\U$, it contains vectors meeting infinitely many
pieces.

\begin{conjecture}[Rigidity]\label{conj:rigidity}
For non-selective $\U$ and the witness subspace $W$ of Theorem~\ref{thm:C}, no
addable blocker exists: every closed $R$ with $R\cap W=0$ satisfies
$\dim(R\cap\HH_S)<\infty$ for some $S\in\U$. Equivalently, $W$ belongs to
\emph{every} pluf extension of $\Phi(\U)$.
\end{conjecture}

\begin{remark}
The conjecture, if true, exhibits a \emph{forced member}: a subspace outside a filter
that nevertheless belongs to every maximalization of it. Nothing of the kind occurs in
a Boolean algebra, where the extension space of a non-maximal filter always splits over
any undecided set. In $\PP(\HH)$ maximalization would then not be the free completion
it is classically, but canalized through transverse subspaces determined by the
combinatorics of $\U$. Whether unique maximalization characterizes selectivity the
conjecture leaves open, since extensions containing $W$ may still differ elsewhere. An
earlier version of this paper asserted incompatible maximalizations for every
non-selective $\U$; that assertion is withdrawn.
\end{remark}

\begin{remark}[Formal verification]\label{rem:lean2}
Sections~\ref{sec:dichotomy}--\ref{sec:zfc} have been formalized and
machine-checked in Lean~4 with Mathlib (toolchain v4.28.0):
Theorem~\ref{thm:A} and
Corollary~\ref{cor:B} (the gliding-hump recursion and its relativization,
the latter proved by re-running the recursion inside a block rather than by
transport along an enumeration); the density clause of
Proposition~\ref{prop:shape}; Lemma~\ref{lem:criteria}, both criteria, and
properness of $\Phi(\U)$; Theorem~\ref{thm:C} in both directions --- necessity
in full, sufficiency parametrized by the conclusion of Mathias's theorem for
the subspace in question, which is the quarantined classical import;
Lemma~\ref{lem:thin}, Proposition~\ref{prop:noblock}, and the construction of
$W$ and the blocks; Theorem~\ref{thm:gowers} in all three clauses (intimacy,
the finite-block dimension bound, and the absence of finite relative
codimension); and Theorem~\ref{thm:D} together with
Corollary~\ref{cor:diag-facts}$(a)$, stated against a named package of the
lattice facts consumed from \cite{FWI} --- upward closure, closure under
meets, properness, membership of the finite-codimension subspaces, and the
maximality criterion --- so that the mechanized statements are theorems of
ZFC. The artifacts compile with no unproved obligations, and every
statement's axiom footprint is exactly \texttt{propext},
\texttt{Classical.choice}, \texttt{Quot.sound}. Three by-products of the
mechanization are worth recording: in Proposition~\ref{prop:noblock} neither
the nonemptiness of the pieces nor their covering $\N$ is used; the
finite-block bound of Theorem~\ref{thm:gowers} admits the sharper proof of
Remark~\ref{rem:pairsum} below; and the extension formulation of
Corollary~\ref{cor:diag-facts}$(a)$ requires its quantifier restricted to
closed subspaces, since the filter generated by the blocks of $\U$ and the
finite-codimension subspaces contains subspaces that are not closed. Theorem~\ref{thm:E} is verified with the continuum hypothesis taken as an
equinumerosity hypothesis, by a route this paper itself makes available:
rather than running the three-case recursion of \cite{FWI}, the mechanization
builds the directed family and appeals to
Proposition~\ref{prop:reduction}, so that maximality, nonprincipality and
non-diagonalizability all follow from results already proved; the blocking
subspace at each stage is then a line, and the three cases collapse to one.
Proposition~\ref{prop:reduction} is verified as printed, quantified over all
orthonormal bases indexed by $\N$ --- which required the basis-relative form
of Theorem~\ref{thm:D}, itself established at an arbitrary Hilbert basis of an
arbitrary real Hilbert space --- and Proposition~\ref{prop:chains} likewise,
its coordinatization resting on a countable orthogonal decomposition adapted to
the chain. The verification artifacts will be archived with the final version
of this paper.
\end{remark}

\begin{remark}
Since selective ultrafilters need not exist (they are absent from the random real model
of Kunen), it is consistent that \emph{no} basis-generated pluf exists, although plufs
always exist by Zorn's lemma. This bears on the strength of the existence of
nonprincipal plufs: the basis route cannot establish it from the Boolean Prime Ideal
theorem, since selectivity exceeds BPI. Whether BPI implies the existence of
nonprincipal plufs we leave open; the natural compactness reduction fails because
maximality in $\PP(\HH)$ (unlike in a Boolean algebra) is an infinitary disjunction over
the blockers of each subspace.
\end{remark}

\section{Intimate subspaces}\label{sec:intimate}

\begin{definition}\label{def:intimate}
A closed subspace $W\subseteq\HH$ is \emph{intimate} (for the basis $(e_n)$) if for
every $A\subseteq\N$, at least one of $W\cap\HH_A$, $W\cap\HH_{\N\setminus A}$ is
nonzero.
\end{definition}

Equivalently: $W$ is non-intimate iff for some two-coloring of the basis, the
coordinate projections of both colors are injective on $W$ --- iff $W$ is, in the
incidence sense, a \emph{bi-graph} over some splitting. Intimacy is thus a partition
property of $W$ against the basis. Trivial examples: any $W$ containing a subspace of
finite codimension in an infinite block $\HH_B$ is intimate, since for any coloring one
color meets $B$ infinitely and finite codimension survives to that side. The trivial examples do not exhaust
intimacy. Gowers's construction follows, with one typographical repair (the display of
the $2$-to-$1$ map) and a concrete choice of the auxiliary vectors.

\begin{theorem}[Gowers, personal communication, c.~2000]\label{thm:gowers}
There is an intimate closed subspace $X$ such that $\dim(X\cap\HH_C)\le|C|/2$ for every
finite $C\subseteq\N$. Consequently $X$ contains no subspace of finite codimension in
any infinite block; the trivial examples do not exhaust intimacy.
\end{theorem}

\begin{proof}
Let $w=(1,\tfrac12,\tfrac13,\dots)\in\ell^2$ and choose vectors $u_n\in\ell^2$ with
first nonzero coordinate at $n$ and $\langle u_n,w\rangle=0$; concretely,
\[
u_n=e_n-\tfrac{n+1}{n}\,e_{n+1}
\qquad\Bigl(\langle u_n,w\rangle=\tfrac1n-\tfrac{n+1}{n}\cdot\tfrac1{n+1}=0\Bigr).
\]
Duplicate: $v_{n,2i-1}=v_{n,2i}=u_{n,i}$, and set $X=\{v_n:n\}^\perp$. Note
$e_{2i-1}-e_{2i}\in X$ for every $i$.

\emph{Intimacy.} Let $\N=A\sqcup B$. If some pair $\{2i-1,2i\}$ lies in $A$, then
$e_{2i-1}-e_{2i}\in X\cap\HH_A$; similarly for $B$. Otherwise every pair splits, so
$A=\{a_1<a_2<\cdots\}$ with $a_i\in\{2i-1,2i\}$; the vector $x=\sum_iw_ie_{a_i}$ lies in
$\ell^2$ and satisfies
$\langle x,v_n\rangle=\sum_iw_iv_{n,a_i}=\sum_iw_iu_{n,i}=\langle w,u_n\rangle=0$, so
$0\neq x\in X\cap\HH_A$ (and symmetrically $X\cap\HH_B\neq0$: in the split case $X$
meets \emph{both} sides).

\emph{The dimension bound.} Let $|C|=n$ and set $I=\{\lceil j/2\rceil:j\in C\}$, so
$|I|\ge n/2$ (the map $j\mapsto\lceil j/2\rceil$ is at worst $2$-to-$1$). For $i\in I$,
the restriction $v_i|_C$ is nonzero: a witness $j\in C\cap\{2i-1,2i\}$ has
$v_{i,j}=u_{i,i}\neq0$. Order $I=\{i_1<i_2<\cdots\}$; since $v_{i}$ vanishes at
positions $<2i-1$, the row $v_{i_{k'}}|_C$ vanishes at the witness position $j_k\le2i_k$
of every earlier row ($2i_k<2i_{k'}-1$): the rows are in echelon form, hence linearly
independent, imposing at least $n/2$ independent conditions on $\HH_C$. So
$\dim(X\cap\HH_C)\le n/2$.

Finally, finite codimension $d$ inside an infinite block $\HH_S$ would force
$\dim(X\cap\HH_C)\ge|C|-d$ for finite $C\subseteq S$, contradicting the bound for
$|C|>2d$.
\end{proof}

\begin{remark}[The finitely supported part of $X$]\label{rem:pairsum}
The dimension bound admits a sharper argument, which identifies the intersection
exactly. Membership in $X$ is the recursion $s_n=\tfrac{n+2}{n+1}s_{n+1}$
on the pair sums $s_n=x_{2n}+x_{2n+1}$ (0-indexed), whose coefficients never
vanish; so if $x\in X$ has finite support, running the recursion downward from a
vanishing tail forces \emph{every} pair sum to vanish. Thus
$X\cap\HH_C=\{x:x_{2k}=-x_{2k+1}\}$ on complete pairs and $x=0$ elsewhere: the
finitely supported part of $X$ is exactly the span of the pair differences
$e_{2k}-e_{2k+1}$, and $\dim(X\cap\HH_C)$ equals the number of complete pairs of
$C$, whence $2\dim(X\cap\HH_C)\le|C|$, with equality iff $C$ is a union of pairs.
\end{remark}

\begin{remark}[The example is finely tuned to its basis]\label{rem:tuned}
With the concrete $u_n$ above, $\{u_n\}^\perp=\R w$ exactly (solve the recursion
$x_n=\tfrac{n+1}{n}x_{n+1}$). Hence for the evens/odds splitting --- where every pair
splits --- the intersections $X\cap\HH_{\mathrm{odd}}$ and $X\cap\HH_{\mathrm{even}}$
are each \emph{one-dimensional}, spanned by the twisted copies of $w$. So $X$ is
intimate only in the nonvanishing sense of Definition~\ref{def:intimate}; its
``infinite-dimensional intimacy'' fails, and $X$ is far from addable to block filters.
For the applications below, nonvanishing intimacy is exactly the right notion.
\end{remark}

Corollary~\ref{cor:B} explains why any counterexample had to look like this.

\begin{proposition}[The shape of nontrivial intimacy]\label{prop:shape}
$W$ is intimate iff $\ZW_W$ contains no complementary pair $\{A,\N\setminus A\}$. If in
addition $W$ has infinite codimension in every block it nearly fills --- i.e.\
$\mathcal C_W=\emptyset$ --- then $\ZW_W$ is dense: every infinite set contains an
infinite $S$ with $W\cap\HH_S=0$. Gowers's $X$ satisfies both conditions.
\end{proposition}

\begin{proof}
The first sentence is the definition. The second is Corollary~\ref{cor:B}. For $X$:
$\mathcal C_X=\emptyset$ by the dimension bound, as in the last step of
Theorem~\ref{thm:gowers}.
\end{proof}

Thus the ``hard case'' of intimacy is precisely: a dense zero-set avoiding complementary
pairs. This is a purely combinatorial specification, and it also explains why such $W$
never obstruct Theorem~\ref{thm:C}: density of $\ZW_W$ is exactly what Mathias's theorem
consumes. Whether the notion admits a structure theory in the style of the Gowers
dichotomies --- via the exact Ramsey principle of Rosendal \cite{Rosendal} and the
Gowers-space framework of de~Rancourt \cite{deRancourt}, in which the relevant pigeonhole
principle visibly fails and the truth should be game-theoretic --- we leave as a program;
the next section shows the pluf application does not wait on it.

\section{Diagonalizable plufs}\label{sec:diag}

\begin{definition}
A pluf $\pluf$ is \emph{diagonalizable} via an orthonormal basis $\mathcal B$ if
$\U^{\mathcal B}_\pluf:=\{S:\HH^{\mathcal B}_S\in\pluf\}$ is an ultrafilter on the index
set of $\mathcal B$. ($\pluf$ is \emph{basic} if $\pluf=\Phi(\U)$ for some basis and
ultrafilter.)
\end{definition}

By \cite{FWI}, Lemma 2.1, a filter $\pluf$ is maximal iff every subspace is a member or
meets some member trivially; and every nonprincipal pluf contains every
finite-codimension subspace. We use both freely.

\begin{theorem}\label{thm:D}
A nonprincipal pluf $\pluf$ is diagonalizable via $\mathcal B$ if and only if every
member of $\pluf$ is $\mathcal B$-intimate.
\end{theorem}

\begin{proof}
($\Rightarrow$) If $R\in\pluf$ satisfies $R\cap\HH_A=R\cap\HH_{A^c}=0$ then neither
$\HH_A$ nor $\HH_{A^c}$ can join $\pluf$, so the trace is not an ultrafilter.

($\Leftarrow$) Given $A$, suppose $\HH_A,\HH_{A^c}\notin\pluf$. By maximality pick
$R_1,R_2\in\pluf$ with $R_1\cap\HH_A=0$ and $R_2\cap\HH_{A^c}=0$; then
$R_1\cap R_2\in\pluf$ is not $\mathcal B$-intimate. So the trace decides every $A$; it
never contains both $A$ and $A^c$ (their blocks meet trivially); and it handles finite
$A$ correctly, since $\HH_{A^c}$ then has finite codimension and belongs to $\pluf$.
Filter properties of the trace follow from those of $\pluf$ via
$\HH_S\cap\HH_{S'}=\HH_{S\cap S'}$.
\end{proof}

\begin{remark}\label{rem:Dgeneral}
The proof of Theorem~\ref{thm:D} uses only the lattice identities
$\HH_S\cap\HH_{S'}=\HH_{S\cap S'}$, monotonicity of $S\mapsto\HH_S$, and the triviality
of $\HH_A\cap\HH_{A^c}$, together with membership of the finite-codimension subspaces.
It therefore holds verbatim for an arbitrary orthonormal basis of an arbitrary real
Hilbert space, with the index set of the basis in place of $\N$; separability plays no
role.
\end{remark}

\begin{corollary}\label{cor:diag-facts}
$(a)$ $\pluf$ is diagonalizable iff $\pluf\supseteq\Phi(\U)$ for some basis and
ultrafilter; in particular diagonalizable nonprincipal plufs exist in ZFC.
$(b)$ $\pluf$ is basic iff it is diagonalizable via a selective ultrafilter
(Theorem~\ref{thm:C}).
$(c)$ Every diagonalizable pluf $\pluf$ has $S_\pluf$ a singleton --- the state-space
face of \cite{FWI}, \S3 --- by Marcus--Spielman--Srivastava; hence the round-slice
property for every ellipsoid. Consequently the ample pluf of \cite{FWI}, Theorem 5.6 is
non-diagonalizable, and the lattice analogue of Anderson's conjecture \cite{Anderson81}
fails under CH.
\end{corollary}

Non-intimacy of a single member survives every enlargement of a filter, so
non-diagonalizability is decided by any family the pluf contains.

\begin{proposition}[One-witness reduction]\label{prop:reduction}
Let $\mathcal R$ be a downward directed family of closed subspaces, all finite
intersections infinite-dimensional, with $\bigcap\mathcal R=0$, such that every
orthonormal basis $\mathcal B$ admits some $R\in\mathcal R$ that is not
$\mathcal B$-intimate. Then every pluf extending $\mathcal R$ (such exist, by Zorn) is
nonprincipal and non-diagonalizable.
\end{proposition}

\begin{proof}
Nonprincipality: $\bigcap\mathcal R=0$ rules out $\bigcap\pluf\neq0$ by \cite{FWI},
Corollary 2.3. Non-diagonalizability via any $\mathcal B$: the witness $R\in\mathcal
R\subseteq\pluf$ and Theorem~\ref{thm:D}.
\end{proof}

Under CH such a family can be built one basis at a time. The construction needs no
spectral theory, and no maximalization is carried through the recursion: only the
family is built, and Proposition~\ref{prop:reduction} supplies the pluf.

\begin{theorem}\label{thm:E}
Assume CH. There is a nonprincipal pluf $\pluf$ that is diagonalizable via no
orthonormal basis.
\end{theorem}

\begin{proof}
Enumerate the orthonormal bases as $\{\mathcal B_i:i<\omega_1\}$ and the nonzero
vectors of $\HH$ as $\{v_i:i<\omega_1\}$; there are $2^{\aleph_0}=\aleph_1$ of each.
We construct closed subspaces $R_i$, $i<\omega_1$, so that, writing $G_i$ for the family
of finite intersections of $\{R_j:j<i\}$ (with $G_0=\{\HH\}$):
\begin{enumerate}
\item[(i)] $\dim(R_i\cap g)=\infty$ for every $g\in G_i$;
\item[(ii)] $R_i$ is not $\mathcal B_i$-intimate;
\item[(iii)] $R_i\cap\R v_i=0$.
\end{enumerate}

Suppose $R_j$ is given for $j<i$. By (i) at earlier stages, $G_i$ is a countable
downward directed family of infinite-dimensional subspaces, so it has a decreasing
cofinal chain $h_1\supseteq h_2\supseteq\cdots$. Write $N=\R v_i$ and
$\mathcal B_i=(e^{(i)}_\beta)$. Choose unit vectors $w_n\in h_n$ and thresholds
$N_1<N_2<\cdots$ recursively so that:
\begin{enumerate}
\item[(a)] $w_n\perp e^{(i)}_\beta$ for all $\beta\le N_n$;
\item[(b)] $w_n\perp w_j$ and $w_n\perp P_{N^\perp}w_j$ for $j<n$, and $w_n\notin N$;
\item[(c)] $w_n$ has at least two nonzero $\mathcal B_i$-coordinates; fix two,
$a_n<b_n$, and take $N_{n+1}>b_n$.
\end{enumerate}
Every condition but the two nonvanishing ones is a finite set of linear constraints,
carving a finite-codimension subspace $C_n$ out of the infinite-dimensional $h_n$; so
$C_n$ is infinite-dimensional. Since $N$ is a line, $C_n\not\subseteq N$, which gives
$w_n\notin N$. And $C_n$ contains a vector with two nonzero coordinates, since a
subspace all of whose vectors are multiples of single basis vectors has dimension at
most one: given independent $u,v\in C_n$, one of $u,v,u+v$ will do.

Set $A=\{a_n:n\}$ and $R_i=\overline{\operatorname{span}}\{w_n\}$. The $w_n$ are
orthonormal by (b), so each $x\in R_i$ has an expansion $x=\sum c_nw_n$. If $L$ is
least with $c_L\neq0$ then, since $w_m$ vanishes on coordinates $\le N_m$ and
$a_L,b_L\le N_{L+1}\le N_m$ for $m>L$,
\[
x_{a_L}=c_Lw_L(a_L)\neq0,\qquad x_{b_L}=c_Lw_L(b_L)\neq0 .
\]
As $a_L\in A$ and $b_L\notin A$, no nonzero $x\in R_i$ lies in $\HH_A$ or in
$\HH_{A^c}$: $R_i\cap\HH_A=R_i\cap\HH_{A^c}=0$, which is (ii). The vectors
$y_n=P_{N^\perp}w_n$ are nonzero by (b), and pairwise orthogonal, since for $j<n$
\[
\langle y_n,y_j\rangle=\langle w_n,P_{N^\perp}y_j\rangle=\langle w_n,y_j\rangle=0 ;
\]
so $P_{N^\perp}x=\sum c_ny_n$ vanishes only for $x=0$, giving $R_i\cap N=0$, which is
(iii). Finally $w_n\in h_n\subseteq h_j$ for $n\ge j$, so
$\dim(R_i\cap h_j)=\infty$ for every $j$, and cofinality of the chain in $G_i$ gives
(i) at every later stage.

Let $\mathcal R$ be the family of finite intersections of the $R_i$. It is downward
directed, and all its members are infinite-dimensional by (i). Every nonzero vector of
$\HH$ is some $v_i$, and $R_i\cap\R v_i=0$, so $\bigcap\mathcal R=0$. Every
orthonormal basis is some $\mathcal B_i$, and $R_i\in\mathcal R$ is not
$\mathcal B_i$-intimate. Proposition~\ref{prop:reduction} applies.
\end{proof}

The ZFC question of \cite{FWI} thus localizes: \emph{exhibit, in ZFC, a single directed
family defeating all bases.}

\section{Toward ZFC}\label{sec:zfc}

We record the constraints known to us on any witness family for
Proposition~\ref{prop:reduction}.

\begin{proposition}[Chains never suffice]\label{prop:chains}
Any countable nested family $M_1\supseteq M_2\supseteq\cdots$ of closed subspaces is
simultaneously coordinatizable: there is a basis $\mathcal B$ for which every $M_k$ is a
block. Blocks are intimate for their own basis; hence every member of the family, and
every enlargement of a member, is $\mathcal B$-intimate, and no countable chain (nor any
family of countable cofinality, applied to a cofinal chain) witnesses
Proposition~\ref{prop:reduction}.
\end{proposition}

\begin{proof}
Choose orthonormal bases of $\bigcap_kM_k$, of each difference $M_k\ominus M_{k+1}$, and
of $\HH\ominus M_1$; their union is a basis making each $M_k$ a block. A block $\HH_B$
is intimate: for any splitting $A\sqcup A^c$, one side meets $B$, and
$\HH_B\cap\HH_A=\HH_{B\cap A}$.
\end{proof}

The natural candidate is therefore \emph{wide}: the full continuous chain
$\mathcal R=\{L^2[0,t]:t\in(0,1]\}$ in $\HH=L^2[0,1]$, with $\bigcap\mathcal R=0$. It is
not simultaneously coordinatizable: the projections $P_{[0,t]}$ generate the continuous
masa $L^\infty[0,1]$, which, being maximal abelian, cannot embed in the atomic masa of
any basis. But non-coordinate does not imply non-intimate, and the question stands:

\begin{question}\label{q:lebesgue}
Is there, for every orthonormal basis $\mathcal B$ of $L^2[0,1]$, some $t$ with
$L^2[0,t]$ not $\mathcal B$-intimate? An affirmative answer gives, via
Proposition~\ref{prop:reduction}, a ZFC nonprincipal non-diagonalizable pluf, settling
the lattice Anderson conjecture outright.
\end{question}

Three orienting remarks. First, the state-theoretic route does not shortcut the
question: the non-diagonalizable pure states of Akemann--Weaver \cite{AW} (under CH) and
Koszmider \cite{Koszmider} (in ZFC) are singular, hence non-regular, and their
$1$-sets are quantum filters that are \emph{not} closed under meets --- they yield no
lattice filter (see \cite{BlecherWeaver}, Lemma 4.3 and the surrounding discussion, and
\cite{FWIII} for the positive counterpart: countably additive pure states \emph{do}
yield plufs). Second, at measurable dimension the question changes character: no family
of coordinate subspaces can defeat all bases there, by the exact dichotomy of
\cite{FWIII}, and a singular countably additive state annihilates every separable-range
projection, so the Lebesgue chain becomes invisible; the correct large-cardinal analogue
of Question~\ref{q:lebesgue} concerns $L^2(\{0,1\}^\kappa)$. Third, Gowers's example
cuts the other way (Remark~\ref{rem:tuned}): its intimacy is genuinely basis-relative,
which is what a witness family must exploit --- one wants a family so symmetric that
every basis is ``generic'' for some member, and the rotation-invariance of the Lebesgue
chain is the obvious source of such symmetry.

\subsection*{Acknowledgements}
Our principal debt is to W.~T.~Gowers, for the correspondence of c.~2000 refuting the
first author's conjecture and posing Galvin's question, and for the construction
reproduced as Theorem~\ref{thm:gowers}; he declined co-authorship at the time, but the
reader will see how much of \S\S2 and 4 is his. We thank Fred Galvin for the question
answered by Theorem~\ref{thm:A}. Portions of these results were developed in the course
of an extended dialogue with Anthropic's Claude language model; responsibility for
correctness rests with the authors.

\begin{thebibliography}{99}

\bibitem{AW} C.~Akemann and N.~Weaver, \emph{Not all pure states on $B(H)$ are
diagonalizable}, Proc. Natl. Acad. Sci. USA \textbf{105} (2008), 5313--5314.

\bibitem{Anderson81} J.~Anderson, \emph{A conjecture concerning the pure states of
$B(H)$ and a related theorem}, in: Topics in Modern Operator Theory, Operator Theory:
Adv. Appl. 2, Birkh\"auser, 1981, 27--43.

\bibitem{BlecherWeaver} D.~P.~Blecher and N.~Weaver, \emph{Quantum measurable
cardinals}, arXiv:1607.08505.

\bibitem{deRancourt} N.~de Rancourt, \emph{Ramsey theory without pigeonhole principle
and the adversarial Ramsey principle}, Trans. Amer. Math. Soc. \textbf{373} (2020),
5025--5056.

\bibitem{FWI} D.~Feldman and A.~Wilce, \emph{Ultrafilter limits in the projection
lattice of a Hilbert space}, in preparation.

\bibitem{FWIII} D.~Feldman and A.~Wilce, \emph{Projection-lattice ultrafilters at a
measurable cardinal}, in preparation.

\bibitem{FWIV} D.~Feldman and A.~Wilce, \emph{Exact paving for Fubini products of
normal measures, with an application to maximal projection filters}, in preparation.

\bibitem{Koszmider} P.~Koszmider, \emph{A non-diagonalizable pure state}, Proc. Natl.
Acad. Sci. USA \textbf{117} (2020), 33084--33089.

\bibitem{Mathias} A.~R.~D.~Mathias, \emph{Happy families}, Ann. Math. Logic \textbf{12}
(1977), 59--111.

\bibitem{MSS} A.~Marcus, D.~A.~Spielman and N.~Srivastava, \emph{Interlacing families
II: Mixed characteristic polynomials and the Kadison--Singer problem}, Ann. of Math.
\textbf{182} (2015), 327--350.

\bibitem{Rosendal} C.~Rosendal, \emph{An exact Ramsey principle for block sequences},
Collect. Math. \textbf{61} (2010), 25--36.

\end{thebibliography}

\end{document}
